\documentclass[11pt]{article}
\usepackage[a4paper,margin=1in]{geometry}
\usepackage[T1]{fontenc}
\usepackage[utf8]{inputenc}
\usepackage{lmodern}
\usepackage{amsmath,amssymb}
\usepackage{booktabs}
\usepackage{xcolor}
\usepackage{hyperref}
\hypersetup{
  colorlinks=true,
  linkcolor=blue!50!black,
  urlcolor=blue!50!black,
  citecolor=blue!50!black
}

\newcommand{\Y}{\mathbf{Y}}
\newcommand{\y}{\mathbf{y}}
\newcommand{\M}{\mathbf{M}}
\newcommand{\A}{\mathbf{A}}
\newcommand{\aev}{\mathbf{a}}
\newcommand{\R}{\mathbf{R}}
\newcommand{\W}{\mathbf{W}}
\newcommand{\Sig}{\boldsymbol{\Sigma}}
\newcommand{\I}{\mathbf{I}}
\newcommand{\m}{\mathbf{m}}
\newcommand{\g}{\mathbf{g}}
\newcommand{\DeltaM}{\boldsymbol{\Delta}}

\title{Critical Review of Candidate Improvements to the Spectreasy Unmixing Method}
\author{Prepared for Spectreasy method development}
\date{\today}

\begin{document}
\maketitle

\section*{Executive Summary}

This report critically evaluates a set of proposed additions to the current
Spectreasy unmixing method. The current Spectreasy path is best understood as
an AutoSpectral-style workflow: it builds a reference matrix from single-color
controls (SCCs), derives autofluorescence (AF) basis signatures, selects an AF
signature per event, performs an ordinary least-squares (OLS) refit, optionally
uses SCC-derived spectral variants, and then applies a Spectreasy marker-only
blend to reduce marker distortion caused by AF projection.

The strongest near-term improvement is \textbf{condition-aware adaptive ridge}.
It directly targets unstable OLS fits caused by similar spectra and can be added
without changing the conceptual identity of Spectreasy. The second most useful
idea is \textbf{shrunken covariance or laser-block whitening}, especially for
residual scoring and AF/variant selection. It is scientifically appealing, but
should be introduced cautiously because a poorly estimated covariance matrix can
silence real biological AF structure. \textbf{Robust residual loss} is a useful
incremental improvement for detector outliers. \textbf{Residual-derived AF
dictionary learning} is promising but should remain conservative and optional.
\textbf{Event-level spectral variability models} are interesting but pose a
high overfitting risk. Finally, \textbf{TRU/FDR-style marker activity selection}
is potentially useful, but it is more a new sparse active-set feature than a
direct improvement of the current Spectreasy implementation.

\section{Current Spectreasy Model}

Let an event's measured spectral detector vector be
\begin{equation}
  \y_i \in \mathbb{R}^{D},
\end{equation}
where \(D\) is the number of spectral detectors. The reference matrix is
\begin{equation}
  \M \in \mathbb{R}^{K \times D},
\end{equation}
where each row is a marker or AF spectrum and \(K\) is the number of reference
rows. The usual linear unmixing model is
\begin{equation}
  \y_i \approx \aev_i \M,
\end{equation}
where
\begin{equation}
  \aev_i \in \mathbb{R}^{K}
\end{equation}
contains the per-event abundances.

For a fixed matrix \(\M\), OLS solves
\begin{equation}
  \hat{\aev}_i
  =
  \arg\min_{\aev_i}
  \left\|
    \y_i - \aev_i \M
  \right\|_2^2.
\end{equation}
Equivalently, if \(\M\M^T\) is well-conditioned,
\begin{equation}
  \hat{\aev}_i
  =
  \y_i \M^T \left(\M\M^T\right)^{-1}.
\end{equation}

Spectreasy adds two important pieces around this basic solve. First, it uses an
AutoSpectral-style AF assignment: for each event, it selects an AF row and fits
the marker rows plus that AF row. Second, it applies a marker-only blend that
pulls selected marker coefficients back toward a marker-only baseline when AF
projection is expected to distort those markers.

This means that the main remaining failure modes are:
\begin{itemize}
  \item unstable coefficients when active spectra are nearly collinear;
  \item detector outliers or saturated channels dominating residual scores;
  \item AF basis signatures not covering the true sample-specific AF;
  \item spectral variants that exist in samples but are not captured by SCCs;
  \item residual scoring that assumes detector errors are independent and equally meaningful.
\end{itemize}

\section{1. Shrunken Covariance and Whitened Residuals}

\subsection*{Mathematical Definition}

The current weighted-least-squares family usually assumes diagonal detector
weights:
\begin{equation}
  \hat{\aev}_i
  =
  \arg\min_{\aev_i}
  \left(\y_i - \aev_i\M\right)^T
  \W_i
  \left(\y_i - \aev_i\M\right),
\end{equation}
where \(\W_i\) is diagonal. The proposed generalization is generalized least
squares (GLS):
\begin{equation}
  \hat{\aev}_i
  =
  \arg\min_{\aev_i}
  \left(\y_i - \aev_i\M\right)^T
  \Sig^{-1}
  \left(\y_i - \aev_i\M\right),
\end{equation}
where \(\Sig\) is a detector residual covariance matrix.

Because full covariance matrices can be noisy or unstable, a shrunken estimate
is safer:
\begin{equation}
  \Sig_{\alpha}
  =
  (1-\alpha)\Sig
  +
  \alpha \operatorname{diag}(\Sig),
  \qquad
  0 \leq \alpha \leq 1.
\end{equation}
The inverse covariance can also be viewed as whitening residuals. If
\begin{equation}
  \Sig_{\alpha} = \mathbf{L}\mathbf{L}^T,
\end{equation}
then the GLS objective is equivalent to
\begin{equation}
  \left\|
    \mathbf{L}^{-1}
    \left(\y_i - \aev_i\M\right)^T
  \right\|_2^2.
\end{equation}

\subsection*{Why It Could Improve Spectreasy}

Spectral detectors are not independent. Neighboring detectors within the same
laser path often share optical, electronic, and biological residual structure.
If two detector errors usually move together, treating them as independent can
make the algorithm overreact to a common noise pattern. A covariance-aware score
tells the method: ``this detector combination is a known background direction,
so do not mistake it for evidence of a marker or AF mismatch.''

For Spectreasy, the most natural first use is not necessarily the final
coefficient solve. It is the scoring step:
\begin{equation}
  s_{i,k}
  =
  \left(\y_i - \hat{\aev}_{i,k}\M_k\right)^T
  \Sig_{\alpha}^{-1}
  \left(\y_i - \hat{\aev}_{i,k}\M_k\right),
\end{equation}
where \(k\) indexes candidate AF rows or candidate spectral variants. This would
select the AF or variant whose residual is unusual after accounting for known
detector covariance, rather than simply the one with the smallest raw squared
error.

\subsection*{Critical View}

This is a real conceptual upgrade, but it is also easy to overdo. A full
\(64 \times 64\) covariance matrix estimated from limited controls can be
ill-conditioned. Worse, covariance estimated from unstained cells may contain
real biological AF patterns. If those patterns are whitened too aggressively,
Spectreasy may remove precisely the structure it needs to model.

A safer implementation would start with laser-block covariance:
\begin{equation}
  \Sig_{\alpha}
  =
  \operatorname{blockdiag}
  \left(
    \Sig_{\alpha}^{UV},
    \Sig_{\alpha}^{V},
    \Sig_{\alpha}^{B},
    \Sig_{\alpha}^{YG},
    \Sig_{\alpha}^{R}
  \right).
\end{equation}
This respects strong within-laser correlations while avoiding fragile
cross-laser covariance. I would also use aggressive shrinkage at first, for
example \(\alpha = 0.5\) to \(0.8\), and expose this only as an experimental
option.

\subsection*{Suggested API}

\begin{verbatim}
unmix_samples(
  unmixing_method = "Spectreasy",
  residual_metric = "laser_block_gls",
  covariance_source = "unstained_residuals",
  covariance_shrinkage = 0.6
)
\end{verbatim}

\section{2. Formal Spectral-Variability Modeling}

\subsection*{Mathematical Definition}

The current model assumes one spectrum per marker, or a small SCC-derived
variant library. A formal spectral-variability model allows each marker spectrum
to move:
\begin{equation}
  \m_{j,i}
  =
  \m_j + \DeltaM_{j,i}.
\end{equation}
An unconstrained version would solve
\begin{equation}
  \min_{\aev_i,\DeltaM_i}
  \left\|
    \y_i - \sum_j a_{ij}(\m_j + \DeltaM_{j,i})
  \right\|_2^2
  +
  \lambda_{\Delta}
  \sum_j
  \left\|
    \DeltaM_{j,i}
  \right\|_2^2.
\end{equation}

A safer, lower-dimensional model is block-wise scaling:
\begin{equation}
  \m_{j,\ell,i}
  =
  s_{j,\ell,i}\m_{j,\ell},
\end{equation}
where \(\ell\) is a laser block. For marker \(j\), the event-level spectrum is
the concatenation of laser blocks:
\begin{equation}
  \m_{j,i}
  =
  \left[
    s_{j,UV,i}\m_{j,UV},
    s_{j,V,i}\m_{j,V},
    s_{j,B,i}\m_{j,B},
    s_{j,YG,i}\m_{j,YG},
    s_{j,R,i}\m_{j,R}
  \right].
\end{equation}

\subsection*{Why It Could Improve Spectreasy}

This addresses a real biological and chemical issue: spectra are not perfectly
fixed. Tandem dyes degrade, detector sensitivity drifts, compensation controls
can differ from samples, and sample conditions can alter fluorescence
properties. A fixed row in \(\M\) may be a poor representation of all events.

Spectreasy already has a pragmatic library method: SCC positive events are
clustered into plausible variant spectra, and per-event optimization can choose
among those variants. A more formal variability model could turn this into a
continuum instead of a small catalog.

\subsection*{Critical View}

This is scientifically attractive but dangerous as a default. If every event can
freely change every marker spectrum, the model can explain away almost anything.
That can reduce residuals while worsening biological accuracy. In plain
language: the method may become too good at making the math look happy.

The main risk is stealing signal. A flexible marker spectrum can absorb AF,
spillover, tandem degradation, or even signal from another marker. That makes
the output harder to interpret. I would not implement free event-level
\(\DeltaM\) as a default Spectreasy component.

A safer path is:
\begin{enumerate}
  \item keep the current SCC-derived spectral-variant library;
  \item add sample-level or batch-level variant diagnostics;
  \item allow block-wise scaling only under strong regularization;
  \item require held-out residual improvement and marker-separation checks.
\end{enumerate}

\subsection*{Suggested API}

\begin{verbatim}
unmix_controls(
  unmixing_method = "Spectreasy",
  spectral_variant_model = "library",
  spectral_variant_max_variants = 8
)

unmix_samples(
  unmixing_method = "Spectreasy",
  spectral_variant_refit = "library"
)
\end{verbatim}

An experimental future option could be:
\begin{verbatim}
spectral_variant_refit = "block_scale"
block_scale_lambda = 10
\end{verbatim}

\section{3. Residual-Derived AF Dictionary Learning}

\subsection*{Mathematical Definition}

After an initial unmixing, residuals are
\begin{equation}
  \R = \Y - \hat{\A}\M.
\end{equation}
The positive residual component is
\begin{equation}
  \R_+ = \max(\R,0).
\end{equation}
A dictionary-learning or nonnegative matrix factorization (NMF) model would
approximate
\begin{equation}
  \R_+
  \approx
  \mathbf{H}\mathbf{G},
\end{equation}
where rows of \(\mathbf{G}\) are candidate residual spectra and
\(\mathbf{H}\) contains event loadings.

Candidate AF spectra should be rejected if they are too marker-like:
\begin{equation}
  \max_j
  \operatorname{cosine}(\g_k,\m_j)
  >
  \tau_{\text{marker}}
  \quad \Rightarrow \quad
  \g_k \text{ is rejected}.
\end{equation}

\subsection*{Why It Could Improve Spectreasy}

The current \texttt{estimate\_af = TRUE} machinery already moves in this direction: it
uses stained-sample residuals, clusters AF-like residual shapes, rejects
marker-like candidates, and scores held-out performance. A more formal AF
dictionary step could help when the unstained control does not contain the AF
structure seen in real stained samples.

This is especially relevant for tissue samples, dead-cell subsets, activated
cells, or preparation-specific autofluorescence. In those cases, the unstained
control may be too narrow. Residual-derived AF learning can act like a second
look: after the model explains known markers, any coherent leftover positive
spectrum may be a missing AF component.

\subsection*{Critical View}

This is promising, but it should not be framed as automatic truth discovery. A
residual component can be missing AF, but it can also be:
\begin{itemize}
  \item unmodeled marker spillover;
  \item tandem degradation;
  \item nonspecific binding;
  \item detector drift;
  \item a real biological marker population that the model failed to fit.
\end{itemize}
The marker-similarity filter is necessary but not sufficient. Some real AF
patterns can be marker-like, and some marker errors can be sufficiently
different to pass the AF filter.

I would keep this as an optional refinement and require held-out improvement:
\begin{equation}
  \operatorname{RMSE}_{\text{heldout}}(\M_{\text{refined}})
  <
  \operatorname{RMSE}_{\text{heldout}}(\M_{\text{base}}) - \epsilon.
\end{equation}

\subsection*{Suggested API}

\begin{verbatim}
unmix_samples(
  unmixing_method = "Spectreasy",
  estimate_af = TRUE,
  af_refine_method = "residual_dictionary",
  af_refine_max_new_bands = 5,
  af_refine_reject_marker_like = TRUE
)
\end{verbatim}

\section{4. Probabilistic Marker Activity Instead of a Fixed TRU Cutoff}

\subsection*{Mathematical Definition}

If a null distribution for marker \(j\) is estimated from unstained cells, let
\begin{equation}
  \hat{F}_j(a)
\end{equation}
be its empirical cumulative distribution function. A one-sided activity
probability can be defined as
\begin{equation}
  p_{ij}
  =
  1 - \hat{F}_j(a_{ij}).
\end{equation}
Then marker \(j\) is active in event \(i\) if
\begin{equation}
  p_{ij} < \alpha.
\end{equation}
With multiple markers, one can apply a false discovery rate (FDR) correction:
\begin{equation}
  K_i
  =
  \left\{
    j : q_{ij} < \alpha
  \right\},
\end{equation}
where \(q_{ij}\) is the FDR-adjusted value for event \(i\).

\subsection*{Why It Could Improve Spectreasy}

If Spectreasy adds sparse active-set refitting, a probabilistic marker activity
step could reduce arbitrary thresholds. Instead of asking whether a coefficient
is above a fixed raw cutoff, the algorithm asks whether the coefficient is
surprising relative to the marker's null behavior.

This could make active-marker selection more stable across markers with
different background spread.

\subsection*{Critical View}

This is not currently replacing a visible fixed TRU step in the present
Spectreasy path. It is better understood as a future sparse refit mode. It also
depends heavily on a good null distribution. If the unstained control is too
small, mismatched, or biologically different from the sample, the empirical CDF
can give false confidence.

FDR also sounds cleaner than it may be. Marker coefficients are correlated
because they are solved from the same detector vector. Standard Benjamini-Hochberg
correction is useful, but it does not magically solve coefficient dependence.

\subsection*{Suggested API}

\begin{verbatim}
unmix_samples(
  unmixing_method = "Spectreasy",
  marker_selection = "null_fdr",
  marker_selection_alpha = 0.005,
  marker_selection_fdr_method = "BH"
)
\end{verbatim}

\section{5. Robust Residual Loss and Huberized Refits}

\subsection*{Mathematical Definition}

Squared residual loss is
\begin{equation}
  L_2(\aev_i)
  =
  \sum_{d=1}^{D}
  r_{id}^2,
  \qquad
  r_{id} = y_{id} - (\aev_i\M)_d.
\end{equation}
A Huber loss is quadratic for small residuals and linear for large residuals:
\begin{equation}
  \rho_k(r)
  =
  \begin{cases}
    \frac{1}{2}r^2, & |r| \leq k, \\
    k(|r|-\frac{1}{2}k), & |r| > k.
  \end{cases}
\end{equation}
The robust objective becomes
\begin{equation}
  \hat{\aev}_i
  =
  \arg\min_{\aev_i}
  \sum_{d=1}^{D}
  \rho_k
  \left(
    y_{id} - (\aev_i\M)_d
  \right).
\end{equation}

\subsection*{Why It Could Improve Spectreasy}

One bad detector should not dominate an event's AF assignment, variant
selection, or final refit. Saturated channels, rare spikes, acquisition
artifacts, and local detector problems can all create large residuals in one or
a few channels. A robust loss reduces their leverage.

This is a natural addition because the code already has Huber-style RWLS logic
for WLS/RWLS. The main improvement would be to let the Spectreasy path use
robust residuals for AF selection and perhaps for the final selected-AF refit.

\subsection*{Critical View}

This is useful but not transformative. It handles outliers; it does not solve
core spectral ambiguity. If two dyes are truly similar, robust loss will not
separate them. If the AF bank is missing the right AF shape, robust loss may only
make the wrong fit less sensitive to the worst channels.

\subsection*{Suggested API}

\begin{verbatim}
unmix_samples(
  unmixing_method = "Spectreasy",
  residual_loss = "huber",
  huber_k = 1.5
)
\end{verbatim}

\section{6. Condition-Aware Adaptive Ridge}

\subsection*{Mathematical Definition}

OLS can be unstable when the active reference matrix is nearly collinear. For an
event-specific active matrix \(\M_i\), define a condition number:
\begin{equation}
  \kappa_i
  =
  \kappa\left(\M_i\M_i^T\right).
\end{equation}
If \(\kappa_i\) is small, use OLS. If it is large, use ridge:
\begin{equation}
  \hat{\aev}_i
  =
  \y_i\M_i^T
  \left(
    \M_i\M_i^T + \lambda_i\I
  \right)^{-1}.
\end{equation}
The ridge strength can increase with ill-conditioning:
\begin{equation}
  \lambda_i
  =
  \lambda_{\min}
  \left(
    \frac{\kappa_i}{\kappa_0}
  \right)^{\gamma},
  \qquad
  \lambda_{\min}
  \leq
  \lambda_i
  \leq
  \lambda_{\max}.
\end{equation}

\subsection*{Why It Could Improve Spectreasy}

This is the most practical improvement. Spectreasy currently relies on OLS in
its AutoSpectral-style path. When the selected AF row and one or more marker
rows are similar, OLS can generate extreme coefficients. These coefficients may
fit the detector vector numerically but look biologically implausible.

Adaptive ridge gives the method a stabilizer. It leaves easy events alone and
only regularizes events where the active matrix is mathematically dangerous. In
simple terms, it stops the solve from using huge positive and negative
coefficients just to explain tiny residual differences.

\subsection*{Critical View}

The main cost is bias. Ridge shrinks coefficients, so very dim real positives
could be slightly underestimated. But this bias is usually preferable to
catastrophic OLS blow-ups in ill-conditioned cases. The key is to apply ridge
only when needed and to report diagnostics.

\subsection*{Suggested API}

\begin{verbatim}
unmix_samples(
  unmixing_method = "Spectreasy",
  adaptive_ridge = TRUE,
  condition_threshold = 1e6,
  ridge_min = 1e-8,
  ridge_max = 1e-2
)
\end{verbatim}

\section{Recommended Development Order}

\begin{table}[h!]
\centering
\begin{tabular}{@{}llll@{}}
\toprule
Rank & Proposal & Expected value & Risk \\
\midrule
1 & Adaptive ridge & High & Low to medium \\
2 & Laser-block covariance scoring & High & Medium \\
3 & Robust residual loss & Medium & Low \\
4 & Residual AF dictionary refinement & Medium to high & Medium to high \\
5 & Probabilistic marker activity & Medium & Medium \\
6 & Event-level spectral variability & Unclear & High \\
\bottomrule
\end{tabular}
\end{table}

\section*{Bottom Line}

The best next Spectreasy improvement is \textbf{adaptive ridge}. It directly
stabilizes the current OLS-based core without forcing a new method identity. The
next best step is \textbf{covariance-aware residual scoring}, ideally using
laser-block covariance with shrinkage. Robust residual loss is a clean
incremental improvement. Residual-derived AF learning is worth keeping and
expanding, but only with strict safeguards. Event-level spectral variability
should remain experimental because it can reduce residuals while silently
damaging biological interpretability.

\end{document}
